\begin{textsample}{POS Dim 3 – rj – Score 19.00 – rj/rj\_inf\_cam\_exp\_13.txt}  \label{ex:f3_pos_006}
Práticas pedagógicas que auxiliam na garantia dos Direitos de Aprendizagem no Campo de Experiências Escuta, Fala, Pensamento e Imaginação A \textbf{professora} pode favorecer que as \textbf{crianças}, desde os \textbf{bebês}, \textbf{brinquem} com a língua oral e a utilizem em situações criadas nos diferentes contextos em que convivem, como nos momentos em que se dá \textbf{colo}, banho, trocas de \textbf{fraldas}, alimentação e nas experiências propostas, em especial, nas \textbf{brincadeiras}, sendo ela uma referência para elas desenvolverem sua fala.
Ela pode perceber avanços nas tentativas de comunicação dos \textbf{bebês}, conforme observa seus \textbf{balbucios}, \textbf{gestos}, expressões faciais, entonação e modulação da voz e os \textbf{ajuda} a organizar seus pedidos, observações, relatos, memórias etc., para que possam pouco a pouco expressar oralmente seus \textbf{desejos} ou sentimentos.
Dentre outras formas comunicativas, as \textbf{crianças} podem ser gradativamente incentivadas a seguir instruções e responder a solicitações compreendendo seus contextos de significação, a elaborar e transmitir recados para diferentes pessoas, a relatar a outra \textbf{criança} um episódio vivido, a formalizar oralmente instruções específicas, tais como regras de jogos, preparo de um prato culinário, procedimentos para \textbf{manipular} um objeto.
Cada um destes meios de comunicação condiciona um tipo de oralidade, um jeito próprio de falar, e cria oportunidades diversas para a \textbf{professora} da Educação \textbf{Infantil}.
Como parceiro sensível no processo de aproximação da \textbf{criança} com a linguagem escrita, a \textbf{professora} observa, \textbf{acolhe} e estimula a refletir sobre aquela linguagem, certa que aprender algo sobre a linguagem que se usa para ler e escrever é um direito a ser garantido na Educação \textbf{Infantil}, respeitando o ritmo e os interesses de cada \textbf{menino} e menina, sem fazer desse processo um ritual de domínio mecânico do código da escrita.
O processo pedagógico neste campo de experiências busca promover vivências nas quais a linguagem verbal na Educação \textbf{Infantil}, aliada a outras linguagens, não seja um conteúdo a ser tratado de modo descontextualizado das práticas sociais significativas das quais a \textbf{criança} participa.
Ele envolve planejamento, frequência e sistematização em função da compreensão do significado que tal linguagem tem na formação cultural da \textbf{criança} e das possíveis formas de sua apropriação.
A preocupação em respeitar o processo de aquisição da linguagem verbal requer tanto reconhecer a participação ativa da \textbf{criança} no processo de significar o mundo, quanto as mediações do professor no arranjo do contexto e na interação com ela.
Conhecer os usos que os \textbf{meninos} e as \textbf{meninas} cotidianamente fazem da linguagem oral e da linguagem escrita, e incorporar esses usos no planejamento didático e nas situações de aprendizagem a serem propostas, impõem à \textbf{professora} trabalhar ludicamente a linguagem oral, a leitura e a escrita com a \textbf{criança} enquanto objetos de reflexão a respeito de suas propriedades.
Desta forma, as interações que \textbf{criança} e \textbf{professora} estabelecem buscam possibilitar que a \textbf{criança} explore a língua, experimente seus \textbf{sons}, diferencie modos de falar, de escrever, reflita por que se fala do jeito que se fala, e por que se escreve do jeito que se escreve.
Cabe à Educação \textbf{Infantil} favorecer, sob a mediação do/a professor/a, oportunidades para a \textbf{criança} \textbf{brincar} com a linguagem oral e a linguagem escrita, e ampliar seus conhecimentos sobre elas, sendo estimulada a formular hipóteses sobre seu funcionamento, testá -las, e empregar estas linguagens nos contextos em que convive.
No domínio da oralidade, a Educação \textbf{Infantil} tem possibilitado às \textbf{crianças} se apropriarem de diversas formas sociais de comunicação, como as cantigas, as \textbf{brincadeiras} de roda, os jogos cantados, e de formas de comunicação presentes na cultura humana : conversas, informações, reclamações, repreensões, elogios etc.
Isto se inicia pela imersão delas em trocas comunicativas e prossegue conforme os momentos de fala criam situações em que elas necessitam pensar sobre a língua, experimentar sua sonoridade, diferenciar maneiras de falar na situação de modo a comunicar \textbf{desejos}, sentimentos, ideias e pensamentos.
Uma forma muito importante de comunicação oral é a conversa, situação em que os sujeitos têm que narrar, descrever, explicar, relatar, ouvir e argumentar com outros parceiros.
É próprio da nossa cultura, \textbf{conversar}, \textbf{contar} casos, o que torna a conversa uma prática social muito frequente.
Por vezes usamos o telefone para isso ou, mais recentemente, aprendemos a usar as redes sociais para trocar ideias e informes com outras pessoas. \textbf{Conversar} é algo que se aprende \textbf{conversando}, e se inicia desde o berçário na Educação \textbf{Infantil}.
Nela a \textbf{professora} é o parceiro com quem a \textbf{criança} inicia suas primeiras conversas com \textbf{balbucios}, \textbf{gestos} etc. e quem a \textbf{ajuda} a organizar seus \textbf{balbucios} em expressões que podem ser compreendidas por qualquer falante de sua língua, iniciando sua apropriação da linguagem oral para relatar \textbf{brincadeiras} ou fatos do cotidiano.

% matched lemmas: acolher, ajuda, balbucio, bebê, brincadeira, brincar, colo, contar, conversar, criança, desejo, fralda, gesto, infantil, manipular, menino, professora, som
\end{textsample}
